\section{Language Server Protocol for \nr{}}
\index{Language Server Protocol}
\index{LSP}

The Language Server Protocol (LSP) is an open standard which most modern editors and 
integrated development environments (IDEs) use to assist programmers with syntax assist 
and code analysis.

LSP servers are language agnostic and communicate, usually over stdin/stdout, with any client 
in the JSON-RPC (JavaScript Object Notation-Remote Procedure Call) protocol.

The full protocol is published at \\ \url{https://microsoft.github.io/language-server-protocol/}.

The \nr{} LSP server is written in \nr{} and implemented as \\ org.netrexx.lsp.NetRexxLspServer. 

The following arguments are accepted:

\begin{tabular}{ l l }
-\/-onChange=nnnn & wait nnnn milliseconds after typing to run diagnostics \\
-\/-onChange=off & run diagnostics on save only \\
-\/-log & log JSON-RPC messages on stderr \\
-\/-stdio & communicate over stdin/stdout \\
-\/-sock  & communicate over socket (development) \\
\end{tabular}

The LSP communication generally flows from the IDE to the LSP server, where the IDE sends a JSON-RSP request 
to the LSP server, upon which the LSP server replies with the appropriate JSON-RPC response. The LSP server 
and the IDE can also send JSON-RPC notifications upon which no response is expected.

LSP servers do not have to implement the full stack of the protocol.
During initial handshake, the IDE client sends an \emph{initialize} JSON-RPC request with which it announces all
protocol features the IDE supports. 
The LSP server then responds documenting its capabilities. When \emph{initialized}, the IDE will restrict its 
requests to the LSP server capabilites.

The \nr{} language server supports the following LSP features: 

\begin{tabular} { |p{4.5cm}|p{8cm}|  }
\hline
TextDocumentSync & Notifications sent from client to server to signal changes to a source file \\
\hline
DocumentSymbols & Requests sent from client to server to request symbol information of a source file \\
\hline
PublishDiagnostics & Notifications sent from server to client to report warnings and errors in a source file \\
\hline
\end{tabular}

When receiving \emph{textDocument/didChange} and \emph{textDocument/didSave} JSON-RPC requests the \nr{} language server 
compiles the source code (without generating a .class file) and sends possible warnings or errors back 
to the client with a \\ 
\emph{textDocument/publishDiagnostics} notification, detailing line number, column and messages if any.\\
Response on \emph{textDocument/didChange} is somewhat delayed (1.5 second by default), so every keystroke would
not trigger a compilation. The delay can be increased by the --onChange argument, or completely switched off 
- a compilation is then only triggered when saving a source file. 

The IDE also sends a \emph{textDocument/documentSymbol} request when an .nrx file is opened or changed,
 the server's response describes the hierarchy of class(es) and method(s) in the 
source file.

The tools/ide-plugins directory of the package contains plugins for the \nr{} Language Server Protocol support.

\begin{itemize}
\item \textbf{lsp4nrx-\emph{Major.Minor.Patch}.vsix} is the plugin for Microsoft's Visual Studio Code. 
Install from \emph{View\~{}Extensions\~{}three dots\~{}Install from VSIX}.
\item \textbf{lsp4nrx-\emph{Major.Minor.Patch}.zip} is the plugin for JetBrains' IntelliJ IDEA.
\item \textbf{\nr{}-UpdateSite.zip} is the plugin for the Eclipse IDE for Java Developers of the
 Eclipse Foundation.
\end{itemize}

All plugins start the language server when any .nrx file is opened.
The language server is started by command \\
'java -cp \emph{dir}/NetRexxC.jar org.netrexx.lsp.NetRexxLspServer', 
where \emph{dir} is the relative path into the plugin to the self-contained current NetRexxC.jar.
This process is terminated by the plugin when the last .nrx file is closed.

All plugins expose the following settings/preferences:
\begin{itemize}
\item \textbf{onChange} controls the compilation delay after typing. When set, the --onChange 
argument is added to the command. 
\item \textbf{log stderr} displays raw JSON-RPC traffic on stderr. When set, the --log argument
 is added to the command.
\item \textbf{classpath} allows to adjust the classpath to find dependent classes. When set, 
its contents is pre-pended on the java -cp classpath command argument.
\end{itemize}

Any LSP compatible IDE has its own specific manner to interface with the language server.
Here we document how to connect the \nr{} language server to Bare Bones Software's BBedit for macOS.

BBEdit defines a new language as a Codeless Language Module. The definition is described in XML plist format.
The \nr{}.plist can be found in the tools/ide-plugins directory of the package.\\
Place the \nr{}.plist in the special dedicated 'Language Modules' folder, and restart BBedit. Go to the 
BBedit Settings and select the Languages section. Click on the +, define .nrx as Extension and 
associate it to \nr{} from the Languages list.\\
Next select 'custom settings', using + again select \nr{} and in the popup window go to the 'Server' tab.\\
Use '/usr/bin/java' as \textbf{Command} and \\
'-cp \emph{dir}/NetRexxC.jar org.netrexx.lsp.NetRexxLspServer' as \textbf{Arguments}.
Adjust \emph{dir} according to where NetRexxC.jar is placed. 
Language ID will be \emph{netrexx}.\\
The icon 'Ready to start the server' will be green. Ensure \textbf{Enable language server} is selected.

